\section*{Beyond Discrete and Continuous Random Variables}

Random variables are typically classified into two types: discrete and continuous. A discrete random variable takes on finitely or countably many values, while a continuous random variable can take any real number as its value. A discrete random variable can be described by a probability mass function (pmf), whereas a continuous random variable can be specified by a probability density function (pdf). By combining discrete and continuous random variables, we can construct mixed-type random variables that contain both discrete and continuous components.

Another type of random variables is known as singular random variables. These random variables are characterized by the property that their probability distribution is concentrated in a region with zero measure, making it impossible to describe them by a probability density function.

In this chapter, we will provide several examples of singular distributions. In later chapters, we will develop a theory that can handle all of these distributions within a single framework. In the second part of this chapter, we will review the definition and basic properties of the Riemann--Stieltjes integral. By using the cumulative distribution function (cdf) of a random variable, we can compute mathematical expectation using Riemann--Stieltjes integral. This provides a unified definition of expectation for discrete, continuous, mixed-type, and singular random variables.

\subsection*{1.1 Discrete and Continuous Random Variables}

In a first course on probability theory, we learn how to model random phenomena using two basic types of random variables. A random variable $X$ of the first type is known as a discrete random variable. If random variable $X$ takes nonnegative integers as its value, we can specify the distribution of $X$ by a probability mass function
\[
p(i)=\Pr(X=i) \qquad \text{for } i=0,1,2,\dots.
\]
Examples of discrete random variables include:
\begin{itemize}
    \item Bernoulli random variable, which has pmf $p(0)=1-p$ and $p(1)=p$ for some $p\in [0,1]$.
    \item Binomial random variable, whose pmf is given by
    \[
    p(i)=\binom{n}{i}p^i(1-p)^{n-i}
    \]
    for $i=0,1,\dots,n$, where $p\in [0,1]$ and $n$ is a positive integer.
    \item Geometric random variable, with pmf given by
    \[
    p(i)=p(1-p)^i
    \]
    for $i=0,1,2,\dots$, where $p$ is a constant between 0 and 1.
    \item Negative binomial random variable, which has pmf
    \[
    p(i)=\binom{i+r-1}{i}(1-p)^i p^r,
    \]
    for some fixed integer $r>0$ and real number $0<p<1$.
    \item Poisson random variable, which has pmf
    \[
    p(i)=\frac{\lambda^i}{i!}e^{-\lambda}
    \]
    for $i\ge 0$, where $\lambda>0$ is a parameter of the distribution.
\end{itemize}

The second type of random variables is known as continuous random variables. A continuous random variable $X$ is described by probability density function, which is a nonnegative function $f(x)$ with integral equal to $1$. The probability that $X$ falls between $a$ and $b$ is computed by the integral $\int_a^b f(x)\, dx$. The followings are common distributions of continuous type:
\begin{itemize}
    \item Uniform distribution, which is described by pdf
    \[
    f(x)=(b-a)^{-1}1_{[a,b]}(x),
    \]
    where $a<b$ are real numbers and $1$ denotes the indicator function
    \begin{equation}
    1_S(x)\triangleq
    \begin{cases}
    1 & \text{if } x\in S,\\
    0 & \text{if } x\notin S.
    \end{cases}
    \tag{1.1}
    \end{equation}
    \item Gaussian distribution with mean $\mu$ and variance $\sigma^2$, which has pdf
    \[
    f(x)=\frac{1}{\sqrt{2\pi\sigma^2}}e^{-(x-\mu)^2/(2\sigma^2)}
    \]
    for $x\in \mathbb{R}$.
    \item Exponential distribution with parameter $\mu>0$, which is specified by the pdf
    \[
    f(x)=\mu e^{-\mu x}1_{[0,\infty)}(x).
    \]
    \item Cauchy distribution, which has pdf
    \[
    f(x)=\frac{1}{\pi(x^2+1)}
    \]
    for $x\in \mathbb{R}$.
    \item Gamma distribution, whose pdf is given by
    \begin{equation}
    \frac{1}{\Gamma(\alpha)\beta^{\alpha}}x^{\alpha-1}e^{-x/\beta}1_{[0,\infty)}(x),
    \tag{1.2}
    \end{equation}
    where $\alpha>0$ and $\beta>0$ are called the shape parameter and the scale parameter, respectively, and
    \[
    \Gamma(\alpha)\triangleq \int_0^{\infty} x^{\alpha-1}e^{-x}\, dx
    \]
    is the Gamma function.
\end{itemize}

The joint Gaussian distribution is a common probability model for multiple random variables. We can regard the random variables as a random column vector. If $\mu$ is an $n$-dimensional column vector and $K$ is an $n\times n$ positive definite matrix, then we define the $n$-dimensional Gaussian distribution by the joint distribution
\[
f_{\mathbf{X}}(\mathbf{x})
=
\frac{1}{(2\pi)^{n/2}\sqrt{\det(K)}}
\exp\left(
-\frac{1}{2}(\mathbf{x}-\mu)^T K^{-1}(\mathbf{x}-\mu)
\right)
\]
with $\mathbf{x}$ ranging over all vectors in $\mathbb{R}^n$. The matrix $K$ is the covariance matrix of the random vector.

When $K$ is the identity matrix, then the random variables in the random vector $\mathbf{X}$ are independent. For example, if $\mu=0$ is the $n$-dimensional zero vector and $K$ is the $n\times n$ identity matrix, the pdf reduces to
\[
f_{\mathbf{X}}(x_1,x_2,\dots,x_n)
=
\frac{1}{(2\pi)^{n/2}}
\exp\left(-(x_1^2+x_2^2+\cdots+x_n^2)/2\right).
\]

This is the joint distribution of $n$ independent standard Gaussian random variables.

If $K$ is positive definite but not equal to the identity, the random variables are correlated. For example, when $n=2$, $\mu=[0,0]^T$, and
\[
K=
\begin{bmatrix}
\sigma_1^2 & \sigma_1\sigma_2\rho\\
\sigma_1\sigma_2\rho & \sigma_2^2
\end{bmatrix}
\]
for some $\sigma_1>0$, $\sigma_2>0$, and $-1<\rho<1$, the joint pdf becomes
\begin{equation}
f_{\mathbf{X}}(x_1,x_2)
=
\frac{1}{2\pi\sigma_1\sigma_2\sqrt{1-\rho^2}}
\exp\left(
-\frac{1}{2(1-\rho^2)}
\left(
\frac{x_1^2}{\sigma_1^2}
+
\frac{x_2^2}{\sigma_2^2}
-
\frac{2\rho x_1y_1}{\sigma_1\sigma_2}
\right)
\right).
\tag{1.3}
\end{equation}

This joint pdf describes a bivariate normal distribution with means $\mu_1=0$ and $\mu_2=0$, variances $\sigma_1^2$ and $\sigma_2^2$, and correlation coefficient $\rho$.

When $\rho=1$ or $\rho=-1$, this pdf is supposed to model perfectly correlated random variables. However, in these cases, we encounter a division-by-zero error in the pdf, and hence, the pdf is not well-defined. In general, when the covariance matrix has a zero determinant, the pdf of the joint Gaussian distribution is not defined. This means that the joint Gaussian distribution is not appropriate for modeling random variables with a singular covariance matrix. In this case, the joint distribution is called a \textit{singular Gaussian distribution}.

For example, consider the case where $X_1$ and $X_2$ have zero mean and the correlation coefficient $\rho$ is equal to $1$. The random variables $X_1$ and $X_2$ are Gaussian distributed individually, but $X_2$ is a deterministic linear function of $X_1$. Strictly speaking, the joint pdf does not exist, but we may use the Dirac delta function $\delta(x)$, which is a kind of generalized functions, to write the joint pdf as
\[
f_X(x_1,x_2)=\delta\left(x_2-\frac{\sigma_2}{\sigma_1}x_1\right)e^{-x_1^2/(2\sigma_1^2)}/\sqrt{2\pi\sigma_1^2}.
\]
